% GENERATED FILE. Edit canonical lessons, then run npm run generate:books.
% canonical-source-hash: fnv1a64:e8d7ca10e37a6624
% canonical-lessons: SA-C399-hradah, SA-C399-snapayati, SA-C399-divyah, SA-C399-toranam, SA-C399-pataka

\chapter{Holy Pools and Gateways}
\label{ch:sa-holy-pools-and-gateways}
\hlchaptermodality{399}

\begin{chapteropening}
\textbf{By the end of this chapter:} \emph{I can describe a holy pool and the gate before a shrine.}

Complete the last lesson of chapter 399: I can describe a holy pool and the gate before a shrine.
\end{chapteropening}

\section[\texorpdfstring{\sk{ह्रदः} (hradaḥ)}{ह्रदः (hradaḥ)}]{\sk{ह्रदः} (hradaḥ) — a deep pool; a lake}

\label{lesson:SA-C399-hradah}



\begin{quote}
Before the new one: say the Sanskrit for purifying, then the Sanskrit for sprinkles over.
\end{quote}

\subsection*{You'll want to know: \sk{ह्रदः}}
\textbf{\sk{ह्रदः}} — \emph{hradaḥ} — \textquotedblleft{}a deep pool; a lake\textquotedblright{}.

\begin{quote}
\sk{वनस्य} \sk{मध्ये} \sk{गभीरः} \sk{ह्रदः} \sk{अस्ति।}
\end{quote}

\begin{quote}
\emph{vanasya madhye gabhīraḥ hradaḥ asti.}
\end{quote}

\begin{quote}
In the middle of the forest there is a deep pool.
\end{quote}

A masculine noun.

\begin{cousinweb}
\textbf{\sk{ह्रदः}} (\emph{hradaḥ}) is the word MW gives for a large or deep piece of water, a lake or pool; MW connects it doubtfully with √hlād, \textquotedblleft{}to refresh\textquotedblright{}, so its origin is uncertain (MW).
\end{cousinweb}

\begin{grammarlens}[title={What you've built}]
One more word for getting around.
\end{grammarlens}

\subsection*{Your turn}
\begin{itemize}
\raggedright
  \item \emph{Say it:} \emph{hradaḥ}
  \item \emph{Say it:} \emph{hradaḥ}, once more
  \item \emph{Say it:} \emph{vanasya madhye gabhīraḥ hradaḥ asti.} — In the middle of the forest there is a deep pool.
  \item \emph{Recall:} say the Sanskrit for purifying, then the Sanskrit for sprinkles over, then say \emph{hradaḥ} again
\end{itemize}

\begin{tcolorbox}[breakable,colback=teal!4,colframe=teal!35!black,title={Before you move on}]
What does \sk{ह्रदः} mean? (\textquotedblleft{}A deep pool; a lake\textquotedblright{}.) Say it once more.
\end{tcolorbox}
\section[\texorpdfstring{\sk{स्नापयति} (snāpayati)}{स्नापयति (snāpayati)}]{\sk{स्नापयति} (snāpayati) — bathes (someone); gives a bath to}

\label{lesson:SA-C399-snapayati}



\begin{quote}
Before the new one: say the Sanskrit for sprinkles over, then the Sanskrit for a deep pool.
\end{quote}

\subsection*{You'll want to know: \sk{स्नापयति}}
\textbf{\sk{स्नापयति}} — \emph{snāpayati} — \textquotedblleft{}bathes (someone); gives a bath to\textquotedblright{}.

\begin{quote}
\sk{माता} \sk{शिशुं} \sk{स्नापयति।}
\end{quote}

\begin{quote}
\emph{mātā śiśuṃ snāpayati.}
\end{quote}

\begin{quote}
Mother gives the baby a bath.
\end{quote}

Past (she …): \sk{सा} \sk{शिशुं} \sk{स्नापितवती}, \emph{sā śiśuṃ snāpitavatī}.

\begin{cousinweb}
\textbf{\sk{स्नापयति}} (\emph{snāpayati}) is built on the root √snā, \textquotedblleft{}to bathe\textquotedblright{}, which gives \emph{snāti}, \textquotedblleft{}bathes\textquotedblright{}; its causative, \textquotedblleft{}makes bathe\textquotedblright{}, is to bathe someone, and MW gives the participle \emph{snāpita} (MW).
\end{cousinweb}

\begin{grammarlens}[title={What you've built}]
One more word for getting around.
\end{grammarlens}

\subsection*{Your turn}
\begin{itemize}
\raggedright
  \item \emph{Say it:} \emph{snāpayati}
  \item \emph{Say it:} \emph{snāpayati}, once more
  \item \emph{Say it:} \emph{mātā śiśuṃ snāpayati.} — Mother gives the baby a bath.
  \item \emph{Recall:} say the Sanskrit for sprinkles over, then the Sanskrit for a deep pool, then say \emph{snāpayati} again
\end{itemize}

\begin{tcolorbox}[breakable,colback=teal!4,colframe=teal!35!black,title={Before you move on}]
What does \sk{स्नापयति} mean? (\textquotedblleft{}Bathes (someone); gives a bath to\textquotedblright{}.) Say it once more.
\end{tcolorbox}
\section[\texorpdfstring{\sk{दिव्यः} (divyaḥ)}{दिव्यः (divyaḥ)}]{\sk{दिव्यः} (divyaḥ) — divine; heavenly}

\label{lesson:SA-C399-divyah}



\begin{quote}
Before the new one: say the Sanskrit for a deep pool, then the Sanskrit for gives a bath to.
\end{quote}

\subsection*{You'll want to know: \sk{दिव्यः}}
\textbf{\sk{दिव्यः}} — \emph{divyaḥ} — \textquotedblleft{}divine; heavenly\textquotedblright{}.

\begin{quote}
\sk{मन्दिरे} \sk{दिव्यः} \sk{प्रकाशः} \sk{अस्ति।}
\end{quote}

\begin{quote}
\emph{mandire divyaḥ prakāśaḥ asti.}
\end{quote}

\begin{quote}
There is a heavenly light in the temple.
\end{quote}

\begin{cousinweb}
\textbf{\sk{दिव्यः}} (\emph{divyaḥ}) is built on \emph{div}, \textquotedblleft{}the sky, heaven\textquotedblright{}, with the ending \emph{-ya}; MW gives \emph{divya} as divine, heavenly, celestial (MW).
\end{cousinweb}

\begin{grammarlens}[title={What you've built}]
One more word for getting around.
\end{grammarlens}

\subsection*{Your turn}
\begin{itemize}
\raggedright
  \item \emph{Say it:} \emph{divyaḥ}
  \item \emph{Say it:} \emph{divyaḥ}, once more
  \item \emph{Say it:} \emph{mandire divyaḥ prakāśaḥ asti.} — There is a heavenly light in the temple.
  \item \emph{Recall:} say the Sanskrit for a deep pool, then the Sanskrit for gives a bath to, then say \emph{divyaḥ} again
\end{itemize}

\begin{tcolorbox}[breakable,colback=teal!4,colframe=teal!35!black,title={Before you move on}]
What does \sk{दिव्यः} mean? (\textquotedblleft{}Divine; heavenly\textquotedblright{}.) Say it once more.
\end{tcolorbox}
\section[\texorpdfstring{\sk{तोरणम्} (toraṇam)}{तोरणम् (toraṇam)}]{\sk{तोरणम्} (toraṇam) — an arched gateway}

\label{lesson:SA-C399-toranam}



\begin{quote}
Before the new one: say the Sanskrit for gives a bath to, then the Sanskrit for divine.
\end{quote}

\subsection*{You'll want to know: \sk{तोरणम्}}
\textbf{\sk{तोरणम्}} — \emph{toraṇam} — \textquotedblleft{}an arched gateway\textquotedblright{}.

\begin{quote}
\sk{मन्दिरस्य} \sk{द्वारे} \sk{तोरणम्} \sk{अस्ति।}
\end{quote}

\begin{quote}
\emph{mandirasya dvāre toraṇam asti.}
\end{quote}

\begin{quote}
There is an arched gateway at the temple door.
\end{quote}

A neuter noun.

\begin{cousinweb}
\textbf{\sk{तोरणम्}} (\emph{toraṇam}) is the word MW gives for an arched doorway or portal, often hung with garlands; MW gives no derivation, so its origin is uncertain; it lives on as Sinhala \emph{toraṇa}, \textquotedblleft{}an arched gateway\textquotedblright{} (MW; CDIAL 5976).
\end{cousinweb}

\begin{grammarlens}[title={What you've built}]
One more word for getting around.
\end{grammarlens}

\subsection*{Your turn}
\begin{itemize}
\raggedright
  \item \emph{Say it:} \emph{toraṇam}
  \item \emph{Say it:} \emph{toraṇam}, once more
  \item \emph{Say it:} \emph{mandirasya dvāre toraṇam asti.} — There is an arched gateway at the temple door.
  \item \emph{Recall:} say the Sanskrit for gives a bath to, then the Sanskrit for divine, then say \emph{toraṇam} again
\end{itemize}

\begin{tcolorbox}[breakable,colback=teal!4,colframe=teal!35!black,title={Before you move on}]
What does \sk{तोरणम्} mean? (\textquotedblleft{}An arched gateway\textquotedblright{}.) Say it once more.
\end{tcolorbox}
\section[\texorpdfstring{\sk{पताका} (patākā)}{पताका (patākā)}]{\sk{पताका} (patākā) — a pennant; a banner}

\label{lesson:SA-C399-pataka}



\begin{quote}
Before the new one: say the Sanskrit for divine, then the Sanskrit for an arched gateway.
\end{quote}

\subsection*{You'll want to know: \sk{पताका}}
\textbf{\sk{पताका}} — \emph{patākā} — \textquotedblleft{}a pennant; a banner\textquotedblright{}.

\begin{quote}
\sk{तोरणे} \sk{पताका} \sk{वायुना} \sk{चलति।}
\end{quote}

\begin{quote}
\emph{toraṇe patākā vāyunā calati.}
\end{quote}

\begin{quote}
A pennant flutters in the wind on the gateway.
\end{quote}

A feminine noun.

\begin{cousinweb}
\textbf{\sk{पताका}} (\emph{patākā}) is the word MW gives for a flag, a pennon or a banner; MW gives no derivation, and CDIAL 7726 takes it as probably from a non-Aryan word, so its origin is uncertain (MW; CDIAL 7726).
\end{cousinweb}

\begin{grammarlens}[title={What you've built}]
One more word for getting around.
\end{grammarlens}

\subsection*{Your turn}
\begin{itemize}
\raggedright
  \item \emph{Say it:} \emph{patākā}
  \item \emph{Say it:} \emph{patākā}, once more
  \item \emph{Say it:} \emph{toraṇe patākā vāyunā calati.} — A pennant flutters in the wind on the gateway.
  \item \emph{Recall:} say the Sanskrit for divine, then the Sanskrit for an arched gateway, then say \emph{patākā} again
\end{itemize}

\begin{tcolorbox}[breakable,colback=teal!4,colframe=teal!35!black,title={Before you move on}]
What does \sk{पताका} mean? (\textquotedblleft{}A pennant; a banner\textquotedblright{}.) Say it once more.
\end{tcolorbox}
